% !TEX root = recommenders_v2.tex
\documentclass[12pt]{article}

\usepackage[top=1in, bottom=1in, left=1in, right=1in]{geometry}
\usepackage[T1]{fontenc}
\usepackage{mathptmx}
\usepackage{amsmath, amssymb, amsthm}
\usepackage{enumitem}
\usepackage{setspace}
\usepackage{xcolor}
\usepackage[colorlinks=true, linkcolor=blue, citecolor=blue, urlcolor=blue]{hyperref}

\linespread{1.15}

\title{The effects of AI recommendation on epistemic diversity}
\author{}
\date{}

\begin{document}
\maketitle

I describe a simple model (sections~\ref{sec:bandit}--\ref{sec:airec}), make some guesses about how the model may behave (section~\ref{sec:conj}), and discuss some potentially philosophically interesting consequences of these conjectures (sections~\ref{sec:improve}--\ref{sec:mech}).

\section{Standard two-armed bandit}
\label{sec:bandit}

Start with a standard two-armed bandit model.
There are two arms, $A$ and $B$.
Arm $A$ has mean payoff $p_A$.
Arm $B$ has mean payoff $p_B$.
Assume $p_A > p_B$, so $A$ is the better arm.

There are $n$ agents.
Assume for now that agents do not communicate or observe others' payoffs.
Each agent has a Beta prior over $p_A$ and $p_B$.
Agents condition on their observed payoffs.

Every agent is a myopic Bayesian.
At every stage, they choose whichever arm has the higher posterior expected payoff.

Let $D_i$ be the event that agent $i$ discovers which is the better arm (this needs to be defined rigorously eventually).
The community succeeds ($D$) if at least one agent succeeds, so
\[
  D = D_1 \cup \cdots \cup D_n.
\]

\section{Unaided collective success}
\label{sec:unaided}

I use $P$ for the ``objective'' probabilities determined by the model specifications.
Suppose that for all $i$
\[
  P(D_i) = q = 0.05.
\]
I don't know how plausible this is, but the point is we're not assuming that any agent is very likely to succeed.

Because the agents are independent, the probability that the community succeeds is much higher than the probability of individual success.
With $n = 100$ agents,
\[
  P(D) = 1 - (1-q)^n = 1 - 0.95^{100} \approx 0.994.
\]

\section{AI recommender}
\label{sec:airec}

Now introduce an AI recommender.
At each stage $t$, agent $i$ gives the recommender her private history $h_i^t$ (which arms she has pulled and which payoffs she has received).
By conditioning on that history, the agent also has posterior $\pi_i^t$.
The recommender $\rho$ takes these two things as inputs and returns a recommendation
\[
  R_i^t := \rho(h_i^t, \pi_i^t, Z) \in \{A, B\}.
\]

The variable $Z$ represents the features of the recommender that are common among all agents, e.g.\ the algorithm that $\rho$ uses to make recommendations.
Suppose (here is a big simplification) $Z \in \{G, F\}$, that is, $Z$ can be in two states (a good state or a failure state).
Suppose
\[
  P(Z = G) = 0.99, \qquad P(Z = F) = 0.01.
\]

In the aided regime, every agent remains a myopic Bayesian but conditions on the AI recommendation before deciding which arm to pull.
So, in the aided regime, first an agent conditions on his history to get the posterior $\pi_i^t$.
Then he conditions on $R_i^t$.
Finally, he chooses the arm with the higher $\pi_i^t(\,\cdot \mid R_i^t)$ expected payoff.

\section{Model conjectures}
\label{sec:conj}

I'm guessing that recommenders and recommendation likelihoods can be defined so that, say,
\[
  P(D_i \mid Z = G) = 0.1, \qquad P(D_i \mid Z = F) = 0.
\]

One idea is that in the good state, the recommender is somewhat helpful and often nudges the agent toward the better arm $A$, raising the probability of discovery from $0.05$ to $0.1$.
In the failure condition, the recommender gives misleading $B$ recommendations that prevent success almost surely.
Suppose, for example, that at early stages, when continued sampling of $A$ is important, the recommender often recommends $B$.
If there are enough early $B$ recommendations, the agent won't discover that $A$ is the good arm.

\section{Individuals improve, the community gets worse}
\label{sec:improve}

For all $i$, in the aided regime, the probability that $i$ succeeds is
\[
  (0.99)(0.1) + (0.01)(0) = 0.099.
\]
So for each individual the probability of discovery improves from $0.05$ to $0.099$.

The community calculation is different because $Z$ is common across all agents.
Conditional on $Z = G$, the agents still have independent private histories, and each has success probability $0.1$.
So the probability that at least one of the $100$ agents succeeds is
\[
  1 - (1 - 0.1)^{100} \approx 0.99997.
\]
Conditional on $Z = F$, no one succeeds.

So, the aided community discovery probability is
\[
  (0.99)\bigl[1 - (1 - 0.1)^{100}\bigr] + (0.01)(0) \approx 0.98997.
\]
The unaided community success probability is about $0.994$.
So the AI raises every individual's probability of discovery while lowering the community's probability of discovery.

\section{The mechanism}
\label{sec:mech}

The mechanism is the introduction of dependence from the common AI recommender.
In the unaided regime, some agents get unlucky and abandon $A$, but others continue sampling it.
Their unluckiness is not correlated, so the community benefits from independent variation in payoff histories.

In the aided regime, the agents still have private evidence, but their evidence now includes recommendations generated by the same AI.
In the good condition, this common system helps many agents.
In the rare failure condition, it moves many posteriors in the same wrong direction.
In the aided regime unluckiness becomes correlated.
That is the sense in which this AI recommendation can decrease epistemic diversity.

Finally, note that the model is not a wheel-shaped network with AI as the central node.
The AI does not sample from the arms or generate any payoff history.

\section{Next steps}
\label{sec:next}

\begin{enumerate}[leftmargin=*]
  \item Fill in the details of the model and run simulations to confirm the conjectures and effects in sections~\ref{sec:conj} and~\ref{sec:improve}.
  \item Think about how to make the model more interesting.
    For example, we could introduce history-sharing networks together with AI recommendations, and we could model more sophisticated AI recommenders.
\end{enumerate}

\end{document}
